\documentclass{article}
\usepackage[UTF8]{ctex}
\usepackage{amsmath,amssymb,amsfonts}
\usepackage{enumitem}
\usepackage{geometry}
\geometry{a4paper, margin=1in}
% === 错题本排版优化（不参与题目标识） ===
\setlength{\parindent}{0pt}
\setlength{\parskip}{0.35em plus 0.10em minus 0.08em}
\setlength{\emergencystretch}{3em}
\setlist{itemsep=0.25em, topsep=0.40em, parsep=0pt}
\allowdisplaybreaks[2]
% === 错题本排版优化结束 ===
\begin{document}
\section*{吉林大学2025年数学分析真题}
\section*{76．吉林大学2025年数学分析试卷}

\subsection{题目1}
1． \(\lim _{x \rightarrow \pi} \frac{\sin (m x)}{\sin ((2 n+1) x)}\) ，其中 \(m, n\) 为正整数．

\subsection{题目2}
2． \(\lim _{n \rightarrow \infty}\left(\frac{2}{\pi} \arctan n\right)^{n}\) ．

\subsection{题目3}
3． \(\lim _{x \rightarrow 0} \frac{\sqrt{\frac{1}{4}-x}+\sqrt[3]{\frac{1}{8}+x^{2}}-e^{2 x}}{\sin x}\) ．

\subsection{题目4}
4． \(\lim _{n \rightarrow \infty}\left(\frac{n}{n^{2}+1}+\frac{n}{n^{2}+2^{2}}+\cdots+\frac{n}{n^{2}+n^{2}}\right)\) ．

\subsection{题目5}
5． \(\lim _{n \rightarrow \infty} n \int_{0}^{1} x^{n} \arctan x \mathrm{~d} x\) ．

\subsection{题目6}
6．设 \(u=\frac{x+z}{y+z}\) ，其中 \(z\) 为由方程 \(z e^{z}=x e^{x}+y e^{y}\) 所确定的 \(x, y\) 的函数，求 \(\frac{\partial u}{\partial x}, \frac{\partial u}{\partial y}\) ．

\subsection{题目7}
7．计算重积分 \(\iiint_{D} \sqrt{x^{2}+y^{2}} \mathrm{~d} x \mathrm{~d} y \mathrm{~d} z\) ，其中 \(D\) 是 \(x^{2}+y^{2}=z^{2}\) 与 \(z=1\) 所围的区域．

\subsection{题目8}
8．计算第二型曲面积分

\[
I=\iint_{\sum} x y^{3} \mathrm{~d} z \mathrm{~d} z+2 x^{3} y \mathrm{~d} z \mathrm{~d} x+3 z^{2} \mathrm{~d} x \mathrm{~d} y
\]

其中曲面 \(\sum\) 是雉面 \(z^{2}=x^{2}+y^{2}\) 在 \(z=0\) 和 \(z=2\) 之间的部分，方向取外侧．

二、求级数 \(\sum_{n=1}^{\infty}\left(n^{2}-1\right) x^{n}\) 的和函数．

三。讨论当 \(p>0\) 为何值时，级数 \(\sum_{n=2}^{\infty} \ln \left(1+\frac{(-1)^{n}}{n^{p}}\right)\) 绝对收敛，条件收敛，发散。

四．设函数 \(f(x)\) 在 \([0,+\infty)\) 上连续，且 \(\lim _{x \rightarrow+\infty} f(x)\) 存在．证明：\(f(x)\) 在 \([0,+\infty)\) 上有界，并且最大值和最小值至少有一个能取到。

五．设函数 \(f(x)\) 在 \([0,1]\) 上连续，在 \((0,1)\) 内可导，\(f(0)=0, f(1)=1\) ．证明：存在 \(\xi, \eta \in(0,1), \xi \neq \eta\) ，使得

\[
f^{\prime}(\xi) f^{\prime}(\eta)=1
\]

六．设函数 \(f(\alpha, \beta)=\int_{\Gamma} \frac{x \mathrm{~d} y-y \mathrm{~d} x}{\alpha x^{2}+\beta y^{2}}, \alpha, \beta>0\) ，其中 \(\Gamma\) 为圆周 \(x^{2}+y^{2}=1\) ，取逆时针方向．

\begin{itemize}
  \item[（1）] 求 \(f(\alpha, \beta)\) 的表达式．
  \item[（2）] 讨论广义重积分
\[
\iint_{D} f(\alpha, \beta)^{p} \mathrm{~d} \alpha \mathrm{~d} \beta
\]
的敛散性，其中 \(p \in(-\infty,+\infty)\) ，区域 \(D=\left\{(\alpha, \beta) \in \mathbb{R}^{2} \mid \alpha^{2}+\beta^{2} \leq 1, \alpha, \beta \geq 0\right\}\) 。
\end{itemize}

七．证明：

\begin{itemize}
  \item[（1）] \(f(x)\) 在 \((a, b)\) 上一致连续当且仅当 \(f(x)\) 在 \((a, b)\) 上连续，并且极限 \(\lim _{x \rightarrow a^{+}} f(x)\) 和 \(\lim _{x \rightarrow b^{-}} f(x)\) 都存在．
\end{itemize}

\subsection{题目4}
原题边界“（2）”：如果 \(f(x)\) 在 \((0,1]\) 上可导，并且极限 \(\lim _{x \rightarrow 0^{+}} x^{2 / 3} f^{\prime}(x)\) 存在，则 \(f(x)+\frac{\sin x}{x}\) 在（0 ，1］上一致连续．
\end{document}
